% ============================================================
\chapter{Applications}
\label{ch:applications}
% ============================================================

\begin{intuition}
TDA has found applications across numerous scientific domains.
Its strength lies in its ability to extract \textbf{qualitative}
information (number of holes, connected components, cavities)
from quantitative data, in a manner that is \textbf{robust to noise}
and \textbf{invariant under deformation}.
\end{intuition}

% ============================================================
\section{Shape analysis}
\label{sec:shapes}
% ============================================================

\begin{definition}[Topological shape descriptor]
\index{shape analysis}
Let $\mathcal{S} \subset \R^3$ be a triangulated surface. Several
filtrations can be constructed:
\begin{enumerate}
  \item \textbf{Geodesic filtration}: $f(x) = d_{\mathcal{S}}(x, x_0)$
        for a base point $x_0$.
  \item \textbf{Curvature filtration}: $f(x) = \kappa(x)$ where
        $\kappa$ is the Gaussian curvature.
  \item \textbf{Heat kernel signature}: $f(x) = h_t(x, x)$, the
        diagonal of the heat kernel at time $t$.
\end{enumerate}
The persistence diagram of each filtration yields a topological
descriptor of the shape.
\end{definition}

\begin{example}[3D shape classification]
Consider a database of 3D models (SHREC, ModelNet). For each object:
\begin{enumerate}
  \item Compute the geodesic distance filtration from several base
        points.
  \item Obtain a set of persistence diagrams in dimensions $0$
        and $1$.
  \item Vectorize (persistence images) and classify (SVM).
\end{enumerate}
This approach is naturally invariant under isometry.
\end{example}

\begin{proposition}[Invariance]
Topological descriptors from geodesic filtrations are invariant under
isometry of the surface. Those from curvature filtrations are
invariant under rigid motions.
\end{proposition}

% ============================================================
\section{Protein structure}
\label{sec:proteins}
% ============================================================

\begin{definition}[Topological representation of a protein]
\index{protein}
A protein is modeled as a point cloud in $\R^3$ (the positions of
the $C_\alpha$ atoms along the backbone). The Rips--Vietoris
filtration on this point cloud captures:
\begin{itemize}
  \item $H_0$: chain connectivity.
  \item $H_1$: loops and secondary structures ($\alpha$-helices,
        $\beta$-sheets).
  \item $H_2$: cavities and pockets.
\end{itemize}
\end{definition}

\begin{example}[Binding pocket detection]
Protein binding pockets (sites where ligands bind) correspond to
topological cavities ($H_2$). Persistent generators in $H_2$
identify these pockets: a point $(b, d)$ with large persistence
$d - b$ signals a \textbf{stable} cavity across scales.
\end{example}

\begin{minted}{python}
import numpy as np
from gudhi import RipsComplex

# C-alpha positions (extracted from a PDB file)
coords = np.loadtxt("protein_ca.xyz")

# Rips complex
rips = RipsComplex(points=coords, max_edge_length=12.0)
simplex_tree = rips.create_simplex_tree(max_dimension=3)
dgm = simplex_tree.persistence()

# Filter H2 (cavities)
h2 = [(b, d) for dim, (b, d) in dgm if dim == 2 and d - b > 1.0]
print(f"Number of significant pockets: {len(h2)}")
\end{minted}

% ============================================================
\section{Time series analysis}
\label{sec:time_series}
% ============================================================

\begin{definition}[Takens embedding]
\index{Takens embedding}
Let $(x_t)_{t=0}^{T}$ be a time series. The \textbf{Takens
embedding}\index{Takens} with dimension $d$ and delay $\tau$ is:
\[
  \mathbf{x}_t = (x_t, x_{t+\tau}, x_{t+2\tau}, \ldots,
  x_{t+(d-1)\tau}) \in \R^d.
\]
\end{definition}

\begin{theorem}[Takens, 1981]
If the time series comes from a dynamical system on a manifold
$\mathcal{M}$ of dimension $m$, then for $d > 2m$ generically,
the Takens embedding is an embedding of $\mathcal{M}$ into $\R^d$.
\end{theorem}

\begin{intuition}
The idea is to reconstruct the phase space of the dynamical system
from a single scalar observation. Then, TDA on the reconstructed
point cloud detects the topology of the attractor: a limit cycle
gives $\beta_1 = 1$, a torus gives $\beta_1 = 2, \beta_2 = 1$, etc.
\end{intuition}

\begin{example}[Periodicity detection]
For a periodic series $x_t = \sin(2\pi t / T)$, the Takens embedding
with $d = 2$ traces an ellipse in $\R^2$. The $H_1$ diagram contains
a highly persistent point, confirming periodicity.
\end{example}

\begin{minted}{python}
import numpy as np
from gtda.time_series import SingleTakensEmbedding

# Periodic time series
t = np.linspace(0, 10 * np.pi, 1000)
x = np.sin(t) + 0.1 * np.random.randn(len(t))

# Takens embedding
embedder = SingleTakensEmbedding(
    time_delay=10, dimension=3, parameters_type="fixed"
)
X_embedded = embedder.fit_transform(x)
\end{minted}

% ============================================================
\section{Image analysis}
\label{sec:images}
% ============================================================

\begin{definition}[Sublevel set filtration for images]
\index{sublevel set filtration}
Let $I : \{0, \ldots, W\} \times \{0, \ldots, H\} \to \R$ be a
grayscale image. The \textbf{sublevel set filtration} is:
\[
  I^{-1}(-\infty, t] = \{(i,j) : I(i,j) \leq t\}.
\]
The persistence of this filtration captures dark connected
components ($H_0$) and bright regions surrounded by dark
ones ($H_1$).
\end{definition}

\begin{example}[Topological segmentation]
In medical imaging, tumors appear as regions of different intensity.
Persistent components in $H_0$ of the sublevel set filtration
identify these regions robustly, without arbitrary thresholding.
\end{example}

% ============================================================
\section{Sensor networks}
\label{sec:sensors}
% ============================================================

\begin{definition}[Topological coverage]
\index{sensor network}
A network of $n$ sensors with range $r$ covers a domain
$\Omega \subset \R^2$. The \v{C}ech complex at radius $r$
built on sensor positions allows detection of \textbf{coverage
holes}: $\beta_1 > 0$ signals a hole in the coverage.
\end{definition}

\begin{theorem}[De Silva--Ghrist, 2007]
\index{De Silva--Ghrist}
If the Rips complex $\mathrm{Rips}_{2r}$ of the sensor positions
has trivial reduced homology and the sensors have detection range
$r_s \geq 2r$, then the domain is covered.
\end{theorem}

\begin{remark}
The advantage of this approach is that it does not require exact
sensor positions (only distances), and it works without coordinates.
\end{remark}

% ============================================================
\section{Materials science}
\label{sec:materials}
% ============================================================

\begin{example}[Porous materials analysis]
\index{porous materials}
Porous materials (zeolites, aerogels) are characterized by their
pore structure. TDA on the atomic representation captures:
\begin{itemize}
  \item $H_0$: grains and aggregates.
  \item $H_1$: channels and tunnels.
  \item $H_2$: enclosed cavities.
\end{itemize}
The persistence of these features correlates with macroscopic
properties (permeability, specific surface area).
\end{example}

\begin{example}[Amorphous metallic glasses]
Nakamura et al.\ (2015) used persistence to distinguish metallic
glasses from liquids: the $H_1$ and $H_2$ diagrams show persistent
structures in the glass that are absent in the liquid, without
requiring an order parameter.
\end{example}

% ============================================================
\section{Exercises}
% ============================================================

\begin{exercise}
Compute the $H_0$ and $H_1$ persistence diagrams of a time series
$x_t = \sin(t) + 0.5 \sin(3t)$ after Takens embedding with $d = 3$.
Interpret the results.
\end{exercise}

\begin{exercise}
For a $100 \times 100$ binary image containing three disks and one
annulus, predict the persistence diagram of the sublevel set
filtration, then verify with \texttt{gudhi}.
\end{exercise}

\begin{exercise}[Sensor network]
Randomly generate $50$ sensors in the unit square. Use $H_1$
persistence to detect coverage holes as a function of the
radius $r$. Plot the number of holes as a function of $r$.
\end{exercise}

\begin{exercise}[Project]
Apply TDA to a real-world dataset of your choice (e.g., EEG signals,
financial data, satellite images). Present the chosen filtration,
the obtained diagrams, and their interpretation.
\end{exercise}
